% 上界与上确界。解析数据：A={1-1/n : n=1,2,...}，图中画出 n=1,...,40；
% sup A=1 不属于 A；全部上界组成 [1,+inf)；1.5 与 2 是两个上界；
% 0.9 不是上界，因为 1-1/11≈0.909>0.9。
\begin{tikzpicture}[line width=.6pt,font=\fontsize{8}{10}\selectfont,>=Stealth,x=46mm,y=1mm]
 % 全部上界组成的区间 [1,+inf)
 \fill[FigureGreenFill] (1,0) rectangle (2.18,2.2);
 \draw[FigureGreen,line width=.75pt] (1,0)--(1,2.2);
 \node[FigureGreen,anchor=south east] at (2.18,2.6) {全部上界：$[1,+\infty)$};
 % 实轴
 \draw[->,FigureStroke,line width=.75pt] (-0.1,0)--(2.22,0) node[right,FigureInk]{$\mathbb R$};
 \draw[FigureStroke,line width=.6pt] (0,0)--(0,-1.2);
 \node[FigureInk,anchor=north] at (0,-1.4) {$0$};
 % 集合 A 的元素
 \foreach \n in {1,...,40}{\fill[FigureBlue] ({1-1/\n},0) circle[radius=.65mm];}
 \draw[FigureBlue,line width=.6pt,decorate,decoration={brace,amplitude=1.6mm}] (0,3.2)--(0.975,3.2);
 \node[FigureBlue,anchor=south] at (0.49,5.2) {$A$的元素$1-1/n$};
 % 上确界：不属于 A
 \filldraw[fill=white,draw=FigurePurple,line width=.9pt] (1,0) circle[radius=.95mm];
 \draw[FigurePurple,line width=.6pt] (1,1.4)--(1.06,8.6);
 \node[FigurePurple,anchor=south west] at (1.02,8.4) {最小上界$\sup A=1\notin A$};
 % 两个普通上界
 \foreach \m in {1.5,2}{\draw[FigureGreen,line width=1.1pt] (\m,-1.6)--(\m,1.6);}
 \node[FigureGreen,anchor=north] at (1.5,-1.8) {上界$1.5$};
 \node[FigureGreen,anchor=north] at (2,-1.8) {上界$2$};
 % 比 1 小的数不是上界
 \draw[FigureRed,line width=1.1pt] (0.9,-1.6)--(0.9,1.6);
 \node[FigureRed,anchor=north east] at (0.92,-1.8) {$0.9$不是上界};
\end{tikzpicture}
